% Generated draft: requires mathematical and source review.
\documentclass[11pt]{article}
\usepackage[margin=25mm]{geometry}
\usepackage{fontspec}
\setmainfont{Latin Modern Roman}
\usepackage{amsmath,amssymb,mathtools,unicode-math}
\setmathfont{Latin Modern Math}
\usepackage{enumitem,needspace,xurl,hyperref}
\setlength{\parindent}{0pt}
\setlength{\parskip}{5pt}
\setlength{\emergencystretch}{3em}
\newcommand{\N}{\mathbb{N}}
\newcommand{\Z}{\mathbb{Z}}
\newcommand{\Q}{\mathbb{Q}}
\newcommand{\R}{\mathbb{R}}
\newcommand{\C}{\mathbb{C}}
\newcommand{\F}{\mathbb{F}}
\newcommand{\A}{\mathbb{A}}
\newcommand{\PP}{\mathbb{P}}
\newcommand{\E}{\mathbb{E}}
\newcommand{\eps}{\varepsilon}
\newcommand{\dd}{\,\mathrm{d}}
\begin{document}
\section*{UnsolvedMath ID: 700102}
\textbf{Source problem code:} catalogue entry
\subsection{Definitions and mathematical statement}
A skew field is understood here to mean a division ring \(D\) with \(1\ne 0\), with commutative division rings allowed. Write \(D^{\times}=D\setminus\{0\}\) for its multiplicative group. Interpret \emph{orderable} in the ordered-division-ring sense: there exists a total order \(<\) on \(D\) such that, for all \(a,b,c\in D\),
\[
a<b\implies a+c<b+c,
\qquad
0<a\ \text{and}\ 0<b\implies 0<ab.
\]
The supplied record does not explicitly define this terminology.

Characterize, by necessary and sufficient conditions expressed in terms of the abstract group, all groups \(G\) satisfying
\[
\exists\text{ an orderable division ring }D\quad G\cong D^{\times}.
\]
There is no finiteness or cardinality restriction in the supplied statement. The requested characterization concerns the full multiplicative group of nonzero elements; an order is not part of the given structure on \(G\).
\subsection{Short English statement}
Determine exactly which abstract groups occur as the multiplicative groups of nonzero elements of orderable division rings. Here orderability is interpreted as the existence of a total order compatible with addition and with multiplication of positive elements.
\subsection{Sources}
\begin{itemize}
\item Local UnsolvedMath registry, supplied complete source record: permanent ID \(700102\), exact source code \texttt{catalogue entry}, title \emph{Kourovka Notebook Problem 2.26}. The supplied status is \texttt{open}; this label has not been independently verified.
\item The supplied background identifies \emph{The Kourovka Notebook (Unsolved Problems in Group Theory)}, second issue (1967), Problem 2.26, as reproduced in the twenty-first issue edition (2026), \texttt{arXiv:1401.0300v48}. The registry describes the collection as edited by E. I. Khukhro and V. D. Mazurov. This reference is supplied attribution, not an independently inspected source.
\item The statement begins with the attribution ``(L. Fuchs).'' The background and proposer metadata instead identify A. I. Kokorin, with proposed year 1967. The supplied record does not explain the relationship between these attributions.
\end{itemize}
\end{document}
